\section{Discussion and conclusion}
Relation-aware compilation can reduce work and storage that arise from
prematurely expanding interactions. The performance measurements support this
opportunity for selected generated-relation workloads. Paging addresses a
different constraint: it makes an explicit billion-edge computation executable
on a memory-limited GPU, while exposing substantial host-staging overhead.
These mechanisms benefit from a shared separation of relation semantics,
traversal and storage, but neither makes every graph program efficient.

The present evidence leaves three important questions open. A same-search
materialized-versus-fused ablation is needed to separate traversal gains from
intermediate elimination. Real-data, wide-feature applications and complete
forward/backward steps are needed to establish learning-workload benefits.
Finally, first-call compilation costs and amortization need to be reported
alongside warm execution. Capacity comparisons additionally need a matched
bounded-memory baseline before claiming an advantage over manual chunking.
These missing measurements limit attribution and generalization; the reported
curves do not substitute for them.

The implementation also has path-dependent differentiation and backend coverage.
CUDA paging remains forward-only, output storage must fit the device, and
distributed ownership is fixed. The current results motivate improving source
staging and output assembly, and assessing load-aware partitioning rather than
assuming that a small halo guarantees strong scaling. Tiga provides a concrete
substrate for exploring these choices without changing the local message and
reducer interface.
